% ============================================================================
%  SECCIÓN 4.9: IMPLEMENTACIÓN DE OPERACIONES CRUD
% ============================================================================

\section{Implementaci\'on de operaciones CRUD}

Las operaciones CRUD (Create, Read, Update, Delete) son la base de la
interacci\'on de la aplicaci\'on con los datos. En este avance se implementan
las operaciones de lectura, creaci\'on y eliminaci\'on, mientras que la
operaci\'on de actualizaci\'on se reserva para un avance futuro.

\subsection{Operaci\'on Read: consulta de escenarios}

La operaci\'on de lectura es la m\'as utilizada en la aplicaci\'on. Los
escenarios se obtienen desde Supabase y se almacenan en el cache de
React Query para evitar consultas redundantes:

\begin{lstlisting}
async function fetchScenarios(): Promise<ScenarioWithDetails[]> {
  const { data, error } = await supabase
    .from('scenarios')
    .select(`
      *,
      scenario_images(url, is_primary, storage_path,
        display_order),
      scenario_sports(sports(nombre))
    `)
    .eq('estado', 'activo')
    .order('nombre');

  if (error) {
    throw new Error(error.message);
  }

  return data as ScenarioWithDetails[];
}
\end{lstlisting}

La consulta incluye las relaciones con im\'agenes, deportes y eventos,
utilizando el sistema de relaciones de Supabase PostgREST.

\subsection{Operaci\'on Read: consulta de un escenario}

La consulta individual de un escenario se ejecuta cuando el usuario
navega a la pantalla de detalle, utilizando el identificador de ruta:

\begin{lstlisting}
async function fetchScenarioById(
  id: string
): Promise<ScenarioWithDetails | null> {
  const { data, error } = await supabase
    .from('scenarios')
    .select(`
      *,
      scenario_images(url, is_primary, storage_path,
        display_order),
      scenario_sports(sports(nombre)),
      events(*)
    `)
    .eq('id', id)
    .single();

  if (error) {
    if (error.code === 'PGRST116') return null;
    throw new Error(error.message);
  }

  return data as ScenarioWithDetails;
}
\end{lstlisting}

Si el escenario no existe (\textit{PGRST116}), la funci\'on retorna
\textit{null} y la interfaz muestra un estado vac\'io.

\subsection{Operaci\'on Create: a\~nadir favorito}

La operaci\'on de creaci\'on se utiliza para a\~nadir un escenario a la
lista de favoritos del usuario. La inserci\'on se realiza en la tabla
\textit{favorites} de Supabase:

\begin{lstlisting}
async function addFavorite(
  userId: string,
  scenarioId: string
): Promise<void> {
  const { error } = await supabase
    .from('favorites')
    .insert({
      user_id: userId,
      scenario_id: scenarioId,
    });

  if (error) {
    throw new Error(error.message);
  }
}
\end{lstlisting}

\subsection{Operaci\'on Delete: eliminar favorito}

La operaci\'on de eliminaci\'on se utiliza para quitar un escenario de
la lista de favoritos:

\begin{lstlisting}
async function removeFavorite(
  userId: string,
  scenarioId: string
): Promise<void> {
  const { error } = await supabase
    .from('favorites')
    .delete()
    .eq('user_id', userId)
    .eq('scenario_id', scenarioId);

  if (error) {
    throw new Error(error.message);
  }
}
\end{lstlisting}

\subsection{Operaci\'on Update: actualizaci\'on de escenario (pendiente)}

La operaci\'on de actualizaci\'on de escenarios se encuentra en
desarrollo y estar\'a disponible en un avance futuro. Esta
funcionalidad permitir\'a a los usuarios con rol de administrador o
gestor modificar la informaci\'on de un escenario existente, incluyendo
su nombre, descripci\'on, capacidad y estado.

La operaci\'on se implementar\'a mediante el m\'etodo \textit{update} de
Supabase, con las siguientes restricciones de seguridad:

\begin{lstlisting}
// Pendiente de implementacion
async function updateScenario(
  id: string,
  updates: Partial<Scenario>
): Promise<void> {
  const { error } = await supabase
    .from('scenarios')
    .update(updates)
    .eq('id', id);

  if (error) {
    throw new Error(error.message);
  }
}
\end{lstlisting}

\subsection{Resumen de operaciones CRUD}

\begin{table}[H]
  \centering
  \caption{Operaciones CRUD implementadas}
  \contenidoConFuente{%
    \begin{tablaAPA}
      \begin{tabular}{@{}p{2.0cm}p{3.6cm}p{5.2cm}p{4.8cm}@{}}
        \toprule
        \textbf{Operaci\'on} & \textbf{Entidad} & \textbf{Funci\'on} &
        \textbf{Estado} \\
        \midrule
        Read & Escenarios & \textit{fetchScenarios()} &
        Implementada. \\
        \addlinespace
        Read & Escenario & \textit{fetchScenarioById()} &
        Implementada. \\
        \addlinespace
        Read & Favoritos & \textit{fetchFavorites()} &
        Implementada. \\
        \addlinespace
        Create & Favorito & \textit{addFavorite()} &
        Implementada. \\
        \addlinespace
        Delete & Favorito & \textit{removeFavorite()} &
        Implementada. \\
        \addlinespace
        Update & Escenario & \textit{updateScenario()} &
        Pendiente. \\
        \bottomrule
      \end{tabular}
    \end{tablaAPA}%
  }{Elaboraci\'on propia en base al c\'odigo fuente del proyecto.}
\end{table}

\subsection{Flujo completo de favoritos}

A continuaci\'on se presenta el flujo completo de la funcionalidad de
favoritos, que involucra las operaciones Read, Create y Delete:

\begin{figure}[H]
\centering
  \figuraAPA{fig:flujo-favoritos}{Flujo completo de la funcionalidad de favoritos}{%
  \begin{tikzpicture}[
  node distance=0.6cm,
  nodo/.style={rectangle, draw=black!60, rounded corners, minimum width=2.8cm,
    minimum height=0.6cm, align=center, font=\small},
  flecha/.style={-Stealth, thick}
]
  \node[nodo] (consulta) {Consultar favoritos};
  \node[nodo, right=of consulta] (muestra) {Mostrar lista};
  \node[nodo, right=of muestra] (press) {Presionar boton};
  \node[nodo, below=of press] (toggle) {Toggle favorito};
  \node[nodo, left=of toggle] (invalidar) {Invalidar cache};
  \node[nodo, left=of invalidar] (actualizar) {Actualizar interfaz};
  \draw[flecha] (consulta) -- (muestra);
  \draw[flecha] (muestra) -- (press);
  \draw[flecha] (press) -- (toggle);
  \draw[flecha] (toggle) -- (invalidar);
  \draw[flecha] (invalidar) -- (actualizar);
\end{tikzpicture}
  }{Elaboraci\'on propia.}
\end{figure}
